\documentclass[11pt]{article}

\usepackage[margin=1in]{geometry}
\usepackage[T1]{fontenc}
\usepackage{lmodern}
\usepackage{microtype}
\usepackage{amsmath,amssymb,amsthm,mathtools}
\usepackage{enumitem}
\usepackage[hidelinks]{hyperref}

\newtheorem{theorem}{Theorem}[section]
\newtheorem{proposition}[theorem]{Proposition}
\newtheorem{lemma}[theorem]{Lemma}
\newtheorem{corollary}[theorem]{Corollary}
\theoremstyle{definition}
\newtheorem{definition}[theorem]{Definition}
\newtheorem{problem}[theorem]{Problem}
\theoremstyle{remark}
\newtheorem{remark}[theorem]{Remark}

\newcommand{\C}{C}
\newcommand{\bN}{\beta\mathbb N}
\newcommand{\cc}{\mathfrak c}
\newcommand{\Pp}{\mathcal P}
\newcommand{\M}{\mathcal M}
\newcommand{\Id}{I}
\newcommand{\supp}{\operatorname{supp}}
\newcommand{\w}{\operatorname{w}}
\newcommand{\dens}{\operatorname{dens}}
\newcommand{\ext}{\operatorname{ext}}
\newcommand{\Kop}{\mathcal K}
\newcommand{\hookc}{\stackrel{c}{\hookrightarrow}}

\title{Complemented Copies in $\C(\bN\times\bN)$:\\
Exact Convex Reductions, Positive Closure Principles,\\
and a Separable Obstruction to the Cube-Averaging Route}
\author{Research note}
\date{June 20, 2026}

\begin{document}
\maketitle

\begin{abstract}
We study the problem of Plebanek, Rondo\v{s}, and Sobota asking whether
$\C(\beta\mathbb N\times\beta\mathbb N)$ contains a complemented isomorphic
copy of $\C(K)$ for every separable compact space $K$.  The problem is not
resolved here, and no full solution or counterexample was located in a
literature search through June 20, 2026.  We prove, however, several exact
reductions and closure results that substantially narrow the target.

First, the problem is equivalent to its restriction to separable compact
convex spaces; in fact it is enough to consider the Bauer simplices
$\Pp(K)$ of probability measures on separable compact spaces $K$.  It is
also enough to consider the weak-star compact symmetric dual balls
$B_{\M(K)}$.  These reductions are quantitative and use norm-one
projections.  Second, every separable almost-Milyutin compactum gives a
positive instance, hence so does every separable Dugundji compactum.  Third,
we isolate a concrete separable dyadic test space $K_*=T^{\mathbb N}$ from a
theorem of Blasco and Ivorra: $K_*$ is not almost Milyutin, while the
canonical averaging constants of the finite powers $T^n$ tend to infinity.
Thus no proof based solely on producing an averaging operator from a Cantor
cube can settle the problem.  Finally, we recast the question as a
complementable-universality problem for
$\Kop(\ell_1,\ell_\infty)$ and prove a no-go lemma showing why the most
direct operator-space projection cannot work.
\end{abstract}

\section{The problem and notation}

All compact spaces are Hausdorff and all Banach spaces are real.  Write
\[
        Z=\C(\bN\times\bN).
\]
The source problem is the following.

\begin{problem}[Plebanek--Rondo\v{s}--Sobota \cite{PRS26}]
Does $Z$ contain a complemented isomorphic copy of $\C(K)$ for every
separable compact space $K$?
\end{problem}

For Banach spaces $X,Y$, the notation $X\hookc Y$ means that there are
bounded operators
\[
        U:X\longrightarrow Y,
        \qquad V:Y\longrightarrow X,
        \qquad VU=\Id_X.
\]
Thus $U$ is an isomorphic embedding and $UV$ is a bounded projection onto
$U[X]$.  We use the quantitative factorization constant
\begin{equation}\label{eq:gamma}
 \gamma_Y(X)=\inf\bigl\{\|U\|\,\|V\|:U:X\to Y,\ V:Y\to X,\ VU=\Id_X\bigr\},
\end{equation}
with the convention that the infimum is $+\infty$ if no such factorization
exists.

The principal imported result from \cite{PRS26} is
\begin{equation}\label{eq:cube-in-Z}
 \C(2^\kappa)\hookc Z
 \qquad (1\leq \kappa\leq\cc).
\end{equation}
Indeed, this is the second assertion of Corollary 1.4 there, applied to the
two-point metric compactum.  The same corollary also gives complemented
copies of $\C([0,1]^\kappa)$.

\section{Two elementary closure principles}

We begin by recording the operator calculation behind the original partial
result.

\begin{lemma}[Transitivity]\label{lem:transitivity}
If $X\hookc Y$ and $Y\hookc W$, then $X\hookc W$.  Quantitatively,
\[
       \gamma_W(X)\leq \gamma_Y(X)\,\gamma_W(Y).
\]
\end{lemma}

\begin{proof}
Choose $U_1:X\to Y$, $V_1:Y\to X$, $U_2:Y\to W$, and $V_2:W\to Y$
with $V_iU_i=\Id$.  Then
\[
       (V_1V_2)(U_2U_1)=\Id_X,
\]
and the norm estimate follows.  In projection language, if $A:Y\to W$ has
complemented range with projection $P$ and $Q:Y\to Y$ projects onto a copy
$J[X]$, then
\[
       \Pi=AQA^{-1}P
\]
projects $W$ onto $AJ[X]$.
\end{proof}

The next form is more flexible than an extension-operator argument.

\begin{proposition}[Averaging closure]\label{prop:averaging-closure}
Let $q:L\to K$ be a continuous surjection between compact spaces.  Suppose
there is a bounded operator
\[
       A:\C(L)\longrightarrow\C(K)
       \quad\text{such that}\quad A q^*=\Id_{\C(K)},
       \qquad q^*f=f\circ q.
\]
If $\C(L)\hookc Z$, then $\C(K)\hookc Z$.
\end{proposition}

\begin{proof}
The identity $Aq^*=\Id$ says exactly that $q^*\C(K)$ is complemented in
$\C(L)$, with projection $q^*A$.  Apply Lemma~\ref{lem:transitivity}.
Explicitly, if $U:\C(L)\to Z$ has complemented range and $R:Z\to U[\C(L)]$
is a projection, then
\[
       \Pi=Uq^*AU^{-1}R
\]
projects $Z$ onto $Uq^*[\C(K)]$.
\end{proof}

A bounded operator $A$ as above is called an \emph{averaging operator} for
$q$.  If $K$ is a closed subspace of a compact space $L$ and the restriction
map $r:\C(L)\to\C(K)$ has a bounded right inverse $E$, then $Er$ is a
projection and the same conclusion follows.  Thus the earlier
extension-operator and retract cases are contained in
Proposition~\ref{prop:averaging-closure}.

\section{An exact reduction to separable Bauer simplices}

The following reduction is the main structural observation of this note.
For a compact space $K$, let $\Pp(K)$ be the weak-star compact space of
probability measures on $K$.

\begin{theorem}[Canonical convexification]\label{thm:convexification}
For every compact space $K$, the Banach space $\C(K)$ is isometric to a
norm-one complemented subspace of $\C(\Pp(K))$.

More precisely, define
\[
 \mathcal A_K:\C(K)\longrightarrow\C(\Pp(K)),
 \qquad (\mathcal A_Kf)(\mu)=\int_K f\,d\mu,
\]
and
\[
 \mathcal R_K:\C(\Pp(K))\longrightarrow\C(K),
 \qquad (\mathcal R_KF)(x)=F(\delta_x).
\]
Then $\|\mathcal A_K\|=\|\mathcal R_K\|=1$,
$\mathcal R_K\mathcal A_K=\Id_{\C(K)}$, and
$\mathcal A_K\mathcal R_K$ is a norm-one projection.  Its range is the
space of continuous affine functions on $\Pp(K)$.
\end{theorem}

\begin{proof}
Weak-star continuity of $\mu\mapsto\mu(f)$ shows that $\mathcal A_K$ is
well defined.  Since the Dirac measures belong to $\Pp(K)$,
\[
 \|\mathcal A_Kf\|
 =\sup_{\mu\in\Pp(K)}\left|\int f\,d\mu\right|
 =\sup_{x\in K}|f(x)|=\|f\|.
\]
The Dirac map $\delta:K\to\Pp(K)$ is continuous, so $\mathcal R_K$ is a
contraction.  Moreover,
\[
 (\mathcal R_K\mathcal A_Kf)(x)
 =\int_K f\,d\delta_x=f(x).
\]
Hence $\mathcal A_K\mathcal R_K$ is a norm-one projection onto
$\mathcal A_K[\C(K)]$.

Every function in that range is continuous and affine.  Conversely, if
$F$ is continuous and affine on $\Pp(K)$, then for every finitely supported
probability measure $\mu=\sum_{j=1}^m a_j\delta_{x_j}$,
\[
       F(\mu)=\sum_{j=1}^m a_jF(\delta_{x_j})
             =(\mathcal A_K\mathcal R_KF)(\mu).
\]
Finitely supported probability measures are weak-star dense in $\Pp(K)$,
so continuity gives $F=\mathcal A_K\mathcal R_KF$ everywhere.
\end{proof}

\begin{lemma}\label{lem:PK-separable}
If $K$ is separable compact, then $\Pp(K)$ is separable and
$\w(\Pp(K))\leq\cc$.
\end{lemma}

\begin{proof}
Let $D\subseteq K$ be countable and dense.  The set of probability measures
of the form
\[
       \sum_{j=1}^m q_j\delta_{d_j},
       \qquad d_j\in D,\quad q_j\in\mathbb Q\cap[0,1],
       \quad \sum_{j=1}^m q_j=1,
\]
is countable and weak-star dense in $\Pp(K)$.  To verify density, take a
basic weak-star neighbourhood determined by finitely many functions
$f_1,\dots,f_r$.  The map
$x\mapsto(f_1(x),\dots,f_r(x))$ has compact range in $\mathbb R^r$; its
integral against a probability measure can be approximated by a finite
convex combination of values on the dense set $D$, and the coefficients can
then be made rational.

Since every member of $\C(K)$ is determined by its restriction to $D$, one
has $|\C(K)|\leq|\mathbb R^D|=\cc$.  The weak-star topology of $\Pp(K)$ is
generated by the coordinate maps $\mu\mapsto\mu(f)$ for $f\in\C(K)$.
Consequently $\w(\Pp(K))\leq\cc$.
\end{proof}

Recall that $\Pp(K)$ is a Bauer simplex: its extreme points are precisely
the Dirac measures, the set
$\ext\Pp(K)=\{\delta_x:x\in K\}$ is closed, and each $\mu\in\Pp(K)$ has the
unique representing probability measure $\delta_*\mu$ on this extreme
boundary.

\begin{theorem}[Exact convex reduction]\label{thm:exact-reduction}
The following assertions are equivalent.
\begin{enumerate}[label=\textup{(\roman*)}]
\item $\C(K)\hookc Z$ for every separable compact space $K$.
\item $\C(S)\hookc Z$ for every separable compact convex space $S$.
\item $\C(\Pp(K))\hookc Z$ for every separable compact space $K$.
\item $\C(B_{\M(K)})\hookc Z$ for every separable compact space $K$, where
      $B_{\M(K)}$ is the weak-star compact unit ball of
      $\M(K)=\C(K)^*$.
\end{enumerate}
Thus, if a counterexample exists, one may choose a separable Bauer simplex,
or a separable weak-star compact symmetric convex dual ball, as a
counterexample.
\end{theorem}

\begin{proof}
The implications (i)$\Rightarrow$(ii), (i)$\Rightarrow$(iii), and
(i)$\Rightarrow$(iv) are immediate, because all spaces appearing in
(ii)--(iv) are separable compact when $K$ is separable.  The assertion for
$\Pp(K)$ follows from Lemma~\ref{lem:PK-separable}.  For the dual ball,
weak-star separability follows because the countable set of finite rational
convex combinations of the measures $\pm\delta_d$, $d\in D$, is weak-star
dense in $B_{\M(K)}$; here we use that the weak-star closed convex hull of
$\{\pm\delta_x:x\in K\}$ is the whole dual unit ball.  The same coordinate
argument as in Lemma~\ref{lem:PK-separable} gives weight at most $\cc$.

For (iii)$\Rightarrow$(i), Theorem~\ref{thm:convexification} gives
$\C(K)\hookc\C(\Pp(K))$, and transitivity finishes the proof.

For (iv)$\Rightarrow$(i), define
\[
 J_K:\C(K)\to\C(B_{\M(K)}),\qquad (J_Kf)(\mu)=\mu(f),
\]
and
\[
 S_K:\C(B_{\M(K)})\to\C(K),\qquad (S_KF)(x)=F(\delta_x).
\]
Then $J_K$ is an isometry, $\|S_K\|=1$, and $S_KJ_K=\Id$.  Hence
$\C(K)$ is norm-one complemented in $\C(B_{\M(K)})$, and transitivity
again applies.
\end{proof}

\begin{remark}
Theorem~\ref{thm:exact-reduction} rules out a common simplification: passing
to probability measures does not merely enlarge the positive class; it is
an exact reformulation of the full problem.  The remaining difficulty is
already present for separable Bauer simplices.
\end{remark}

\section{The almost-Milyutin positive class}

A compact space $K$ is called \emph{almost Milyutin} if there exist a
Cantor cube $2^\Gamma$, a continuous surjection $q:2^\Gamma\to K$, and a
bounded averaging operator $A:\C(2^\Gamma)\to\C(K)$ for $q$.  If $A$ can be
chosen regular and positive, one obtains the usual Milyutin property.

\begin{lemma}[Coordinate reduction]\label{lem:coordinate-reduction}
Suppose $q:2^\Gamma\to K$ is a continuous surjection onto a compact space
of weight $\kappa$ and $q$ admits a bounded averaging operator.  Then there
are $\Lambda\subseteq\Gamma$ with $|\Lambda|\leq\kappa$, a continuous
surjection $q_\Lambda:2^\Lambda\to K$, and a bounded averaging operator for
$q_\Lambda$ of no larger norm.
\end{lemma}

\begin{proof}
Embed $K$ into $[0,1]^\kappa$.  Each scalar coordinate of the composite
$2^\Gamma\to K\to[0,1]^\kappa$ depends on only countably many coordinates
of $2^\Gamma$.  Taking the union of these coordinate sets gives
$\Lambda\subseteq\Gamma$ with $|\Lambda|\leq\kappa$ such that
$q=q_\Lambda\circ\pi_\Lambda$ for some onto map
$q_\Lambda:2^\Lambda\to K$.

If $A$ averages $q$, define
\[
       A_\Lambda g=A(g\circ\pi_\Lambda),
       \qquad g\in\C(2^\Lambda).
\]
Then $\|A_\Lambda\|\leq\|A\|$ and
$A_\Lambda q_\Lambda^*=\Id_{\C(K)}$.
\end{proof}

\begin{theorem}\label{thm:almost-milyutin-positive}
If $K$ is a separable almost-Milyutin compact space, then
$\C(K)\hookc Z$.  In particular, the answer to Problem~1 is positive for
every separable Dugundji compact space.
\end{theorem}

\begin{proof}
A separable compact space has weight at most $\cc$.  By
Lemma~\ref{lem:coordinate-reduction}, the almost-Milyutin witness can be
chosen with $|\Gamma|\leq\cc$.  Hence
\[
       \C(K)\hookc\C(2^\Gamma)\hookc Z,
\]
where the first factorization comes from the averaging operator and the
second from \eqref{eq:cube-in-Z}.

Haydon proved that every Dugundji compactum is Milyutin
\cite{Haydon74}; therefore the final assertion follows.
\end{proof}

This theorem strictly extends the elementary retract and extension-operator
subcases, but it still does not cover all separable compact spaces, as the
next section shows.

\section{A separable dyadic obstruction to the cube route}

Blasco and Ivorra constructed a particularly sharp obstruction to bounded
averaging operators from Cantor cubes \cite{BlascoIvorra94}.  We recall the
construction in a form tailored to the present problem.

Let $\alpha=\omega_1$ and let $T$ be the quotient of $2^\alpha$ obtained by
identifying two distinct points.  For a compact space $L$, put
\begin{equation}\label{eq:pL}
 p(L)=\inf\bigl\{\|A\|:\ q:2^\Gamma\twoheadrightarrow L,\
       A:\C(2^\Gamma)\to\C(L),\ Aq^*=\Id_{\C(L)}\bigr\},
\end{equation}
with $p(L)=+\infty$ if no such pair exists.

\begin{theorem}[Blasco--Ivorra]\label{thm:BI}
For every positive integer $n$,
\[
       p(T^n)\geq n+1,
\]
although each finite power $T^n$ is almost Milyutin.  Moreover,
\[
       K_*=T^{\mathbb N}
\]
is dyadic but not almost Milyutin; equivalently, $p(K_*)=+\infty$.
\end{theorem}

The space $K_*$ is a legitimate test object for Problem~1.

\begin{proposition}\label{prop:Kstar-separable}
The compact space $K_*$ is separable, dyadic, and has weight at most
$\cc$.
\end{proposition}

\begin{proof}
The Hewitt--Marczewski--Pondiczery theorem implies that $2^{\omega_1}$ is
separable, since $\omega_1\leq\cc$; see Ross and Stone \cite{RossStone64}.
Thus its quotient $T$ is separable, and a countable product of separable
spaces is separable.  If $q:2^\alpha\to T$ is the quotient map, then
\[
       q^{\mathbb N}:(2^\alpha)^{\mathbb N}\longrightarrow T^{\mathbb N}
\]
is onto, while $(2^\alpha)^{\mathbb N}\cong2^{\alpha\times\mathbb N}
\cong2^\alpha$.  Hence $K_*$ is dyadic.  The weight estimate is immediate.
\end{proof}

Consequently, a proof of Problem~1 cannot proceed by showing that every
separable compact space is almost Milyutin.  The gap between canonical
averaging and arbitrary isomorphic complementability is essential.
The following quantitative observation makes this gap precise.

\begin{proposition}[Uniform finite-stage test]\label{prop:uniform-test}
If $\C(K_*)\hookc Z$, then
\[
       \sup_{n\geq1}\gamma_Z\bigl(\C(T^n)\bigr)<\infty.
\]
More precisely,
\[
       \gamma_Z\bigl(\C(T^n)\bigr)
       \leq \gamma_Z\bigl(\C(K_*)\bigr)
       \qquad(n\geq1).
\]
\end{proposition}

\begin{proof}
Fix $t_0\in T$.  The coordinate projection
$\pi_n:K_*\to T^n$ induces an isometric embedding
\[
 J_n:\C(T^n)\to\C(K_*),\qquad J_nf=f\circ\pi_n.
\]
The map
\[
 R_n:\C(K_*)\to\C(T^n),\qquad
 R_ng(x_1,\dots,x_n)=g(x_1,\dots,x_n,t_0,t_0,\dots)
\]
is a contraction and $R_nJ_n=\Id$.  Composing any factorization of
$\Id_{\C(K_*)}$ through $Z$ with $J_n$ and $R_n$ gives the required
factorization for $\C(T^n)$ with no larger norm product.
\end{proof}

Combining Theorem~\ref{thm:BI} and
Proposition~\ref{prop:uniform-test}, we obtain a concrete possible route to
a negative answer: it would suffice to prove, for the finite powers $T^n$,
a bound of the form
\begin{equation}\label{eq:missing-comparison}
       p(T^n)\leq \Phi\!\left(\gamma_Z(\C(T^n))\right)
\end{equation}
for one finite-valued function $\Phi$ independent of $n$.  Indeed, the
right-hand side would remain bounded if $\C(K_*)\hookc Z$, while the
left-hand side is at least $n+1$.  No estimate such as
\eqref{eq:missing-comparison} is currently available: the Blasco--Ivorra
constant concerns canonical pullback copies arising from continuous maps,
whereas $\gamma_Z$ allows arbitrary isomorphic copies.  This is the precise
point at which the candidate stops short of a counterexample.

\section{The compact-operator reformulation}

There are natural isometric identifications
\begin{equation}\label{eq:operator-identification}
 \begin{split}
 Z
 &\cong \C(\bN,\C(\bN))
  \cong \C(\bN,\ell_\infty)\\
 &\cong \Kop(\ell_1,\ell_\infty)
  \cong \ell_\infty\widehat\otimes_\varepsilon\ell_\infty.
 \end{split}
\end{equation}
For the middle identification, an operator
$S:\ell_1\to\ell_\infty$ is represented by the sequence
$(Se_n)_{n\in\mathbb N}$.  It is compact exactly when this sequence has
relatively norm-compact range; such a sequence extends uniquely to a
norm-continuous map from $\bN$ into $\ell_\infty$.

If $K$ is separable and $(d_n)$ is dense in $K$, then
\begin{equation}\label{eq:dense-evaluation}
       E_K:\C(K)\longrightarrow\ell_\infty,
       \qquad E_Kf=(f(d_n))_{n\in\mathbb N},
\end{equation}
is an isometry.  Thus Problem~1 asks whether
$\Kop(\ell_1,\ell_\infty)$ is complementably universal for those
$\C(K)$-spaces that admit a countable $1$-norming set.

A tempting strategy is to place $\Kop(\ell_1,E_K\C(K))$ inside
$\Kop(\ell_1,\ell_\infty)$ and try to project onto it.  The next lemma shows
that this strategy cannot give anything new.

\begin{lemma}[Operator-space no-go lemma]\label{lem:no-go}
Let $X$ be a closed subspace of a Banach space $Y$.  If
$\Kop(\ell_1,X)$ is complemented in $\Kop(\ell_1,Y)$, then $X$ is
complemented in $Y$.  Moreover, a projection of norm $M$ in the operator
space produces a projection of norm at most $M$ in $Y$.
\end{lemma}

\begin{proof}
For $y\in Y$, define the rank-one operator
\[
       R_y:\ell_1\to Y,\qquad R_y(a)=a_1y.
\]
Suppose $Q:\Kop(\ell_1,Y)\to\Kop(\ell_1,X)$ is a projection.  Set
\[
       P(y)=Q(R_y)e_1.
\]
Then $P:Y\to X$ is bounded with $\|P\|\leq\|Q\|$.  If $x\in X$, then
$R_x\in\Kop(\ell_1,X)$, so $Q(R_x)=R_x$ and $P(x)=x$.  Hence $P$ is a
projection onto $X$.
\end{proof}

Therefore a positive solution must use a genuinely twisted copy of
$\C(K)$ inside $\Kop(\ell_1,\ell_\infty)$; it cannot simply project all
compact operators into the evaluation copy of $\C(K)$ inside
$\ell_\infty$.

\section{A sufficient coupling mechanism}

The source paper succeeds for compact groups by constructing continuous
families of probability couplings supported on the fibres of a difference
map.  The abstract mechanism is short enough to state separately.

\begin{proposition}[Fibre-supported coupling criterion]\label{prop:coupling}
Let $\rho:L\to K$ be a continuous surjection between compact spaces.
Suppose there is a weak-star continuous map
\[
       K\longrightarrow\Pp(L),\qquad x\longmapsto\nu_x,
\]
such that $\supp\nu_x\subseteq\rho^{-1}(x)$ for every $x\in K$.  Then
\[
       Af(x)=\int_L f\,d\nu_x
\]
defines a norm-one positive averaging operator for $\rho$.  Consequently,
if $\C(L)\hookc Z$, then $\C(K)\hookc Z$.
\end{proposition}

\begin{proof}
Weak-star continuity gives $Af\in\C(K)$ and $\|A\|=1$.  For
$g\in\C(K)$,
\[
 A(g\circ\rho)(x)
 =\int_{\rho^{-1}(x)}g(\rho(y))\,d\nu_x(y)=g(x).
\]
Apply Proposition~\ref{prop:averaging-closure}.
\end{proof}

A full positive solution would follow, for example, from a theorem that
every separable compact space is the target of such a map from a compact
$L$ whose $\C(L)$ is already known to be complemented in $Z$.  The
Blasco--Ivorra example shows that one cannot always take $L$ to be a Cantor
cube.  The compact-group construction in \cite{PRS26} supplies exactly this
criterion in a highly structured setting, but no general replacement for
the group operation is presently known.

\section{Status of the literature and conclusion}

The exact wording of Problem~1, the title of \cite{PRS26}, its citation
trail, and combinations of the terms
``$\C(\beta\mathbb N\times\beta\mathbb N)$'', ``separable compact'',
``complemented'', ``Milyutin'', and ``averaging operator'' were searched
through June 20, 2026.  No paper giving a complete affirmative answer or a
counterexample was located.  The recent preprint of Rondo\v{s} and Sobota
\cite{RS26} gives a powerful characterization of complemented copies
$\C(L)$ when $L$ is metrizable, hence when $\C(L)$ is separable as a Banach
space.  It does not cover the nonmetrizable separable compacta that form the
unresolved part of Problem~1.

The strongest rigorous conclusions obtained here are:
\begin{enumerate}[label=\arabic*.]
\item The original question is \emph{exactly equivalent} to the same
      question for separable Bauer simplices $\Pp(K)$, and also to the same
      question for separable symmetric weak-star dual balls
      $B_{\M(K)}$.
\item Every separable almost-Milyutin compactum, in particular every
      separable Dugundji compactum, gives a positive instance.
\item The separable dyadic compactum $K_*=T^{\mathbb N}$ is not almost
      Milyutin and therefore blocks every universal proof based on a
      bounded averaging from a Cantor cube.  Its finite powers provide the
      quantitative test $p(T^n)\geq n+1$.
\item In the operator model
      $Z\cong\Kop(\ell_1,\ell_\infty)$, the obvious attempt to project onto
      $\Kop(\ell_1,\C(K))$ is impossible unless the original evaluation
      copy of $\C(K)$ was already complemented in $\ell_\infty$.
\end{enumerate}

Thus the full problem remains open in this note.  The two sharpest next
targets are either to prove that the Blasco--Ivorra space $K_*$ cannot
factor through $Z$, perhaps by establishing a comparison such as
\eqref{eq:missing-comparison}, or to build a noncanonical twisted
factorization of every separable Bauer simplex through
$\Kop(\ell_1,\ell_\infty)$.

\begin{thebibliography}{99}

\bibitem{BlascoIvorra94}
J.~L. Blasco and C. Ivorra,
\emph{On dyadic spaces and almost Milyutin spaces},
Colloq. Math. \textbf{67} (1994), 241--244.

\bibitem{Ditor73}
S.~Z. Ditor,
\emph{Averaging operators in $\C(S)$ and lower semicontinuous sections of
continuous maps},
Trans. Amer. Math. Soc. \textbf{175} (1973), 195--208.

\bibitem{Haydon74}
R. Haydon,
\emph{On a problem of Pe\l czy\'nski: Milutin spaces, Dugundji spaces and
$AE(0\text{-dim})$},
Studia Math. \textbf{52} (1974), 23--31.

\bibitem{Pelczynski68}
A. Pe\l czy\'nski,
\emph{Linear extensions, linear averagings, and their applications to
linear topological classification of spaces of continuous functions},
Dissertationes Math. (Rozprawy Mat.) \textbf{58} (1968), 1--89.

\bibitem{PRS26}
G. Plebanek, J. Rondo\v{s}, and D. Sobota,
\emph{Complemented subspaces of Banach spaces $\C(K\times L)$},
J. Funct. Anal. \textbf{290} (2026), no.~2, Paper No.~111236, 20 pp.;
arXiv:2405.19120.

\bibitem{RS26}
J. Rondo\v{s} and D. Sobota,
\emph{Complementability of separable spaces $\C(K)$ in Banach spaces},
arXiv:2603.12922, 2026.

\bibitem{RossStone64}
K.~A. Ross and A.~H. Stone,
\emph{Products of separable spaces},
Amer. Math. Monthly \textbf{71} (1964), 398--403.

\end{thebibliography}

\end{document}
